% ============================================================
% APPENDICES
% Place after References or before, depending on JCAP style
% ============================================================

\appendix

% ============================================================
% APPENDIX A — Derivation Details
% ============================================================
\section{Derivation of the comoving IR cutoff}
\label{app:derivation}

This appendix presents the complete derivation of eq.~(\ref{eq:kc_comoving})
in explicit detail.

\paragraph{Step 1: Gibbons--Hawking entropy.}
The de~Sitter horizon has area $A = 4\pi R_{\rm dS}^2 = 4\pi H_{\rm inf}^{-2}$.
The Gibbons--Hawking entropy~\cite{Gibbons:1977} is
\begin{equation}
  S_{\rm dS} = \frac{k_B\,A}{4\,l_P^2}
  = \frac{k_B\,\pi}{H_{\rm inf}^2\,l_P^2}\,.
  \label{eq:app_SdS}
\end{equation}

\paragraph{Step 2: Hilbert space dimension.}
The CCR framework~\cite{Chronis:2026} identifies $S_{\rm dS}/k_B = \ln D$,
giving
\begin{equation}
  \ln D = \frac{\pi}{H_{\rm inf}^2\,l_P^2}
  \qquad\Longrightarrow\qquad
  H_{\rm inf} = \frac{1}{l_P}\sqrt{\frac{\pi}{\ln D}}\,.
  \label{eq:app_Hinf}
\end{equation}

\noindent\textit{Saturation of the holographic bound.}
The use of equality rather than inequality in eq.~(\ref{eq:app_Hinf})
--- i.e.\ $\ln D = S_{\rm dS}/k_B$ rather than $\ln D \leq S_{\rm dS}/k_B$
--- corresponds to the assumption that the de~Sitter horizon entropy
saturates the Bekenstein--Hawking bound, as is standard in de~Sitter
thermodynamics~\cite{Gibbons:1977}.  This is the gravitational
analogue of a black hole whose entropy equals $A/(4\,l_P^2)$: the
horizon is a maximum-entropy surface for the enclosed degrees of
freedom.  If the bound is not saturated ($\ln D < S_{\rm dS}/k_B$),
the effective Hilbert space dimension is smaller, $H_{\rm inf}$ is
larger, and the infrared cutoff $k_c$ shifts to larger values ---
producing \emph{stronger} suppression than the fiducial case.  The
saturation assumption therefore yields the most conservative (i.e.\
weakest) CCR prediction, and any sub-saturation would only strengthen
the observational signature.

For the fiducial value $\ln D = 2 \times 10^{14}$:
\begin{equation}
  H_{\rm inf} = \frac{\sqrt{\pi/(2\times 10^{14})}}{1.616\times 10^{-35}\;\mathrm{m}}
  = \frac{1.253\times 10^{-7}}{1.616\times 10^{-35}\;\mathrm{m}}
  \simeq 7.75\times 10^{27}\;\mathrm{m}^{-1} \;.
\end{equation}
Converting to GeV using $\hbar c = 1.973\times 10^{-16}\;\mathrm{GeV\cdot m}$:
\begin{equation}
  H_{\rm inf} = 7.75\times 10^{27}\;\mathrm{m}^{-1}
  \times 1.973\times 10^{-16}\;\mathrm{GeV\cdot m}
  \simeq 1.53\times 10^{12}\;\mathrm{GeV} \;,
\end{equation}
consistent with table~1.

\paragraph{Step 3: Physical IR cutoff.}
The horizon-crossing condition gives the physical cutoff:
\begin{equation}
  k_{c,\,\rm phys} = H_{\rm inf} = \frac{1}{l_P}\sqrt{\frac{\pi}{\ln D}}\,.
  \label{eq:app_kc_phys}
\end{equation}

\paragraph{Step 4: Scale factor at the onset of inflation.}
Entropy conservation after reheating gives (normalising $a_0 = 1$):
\begin{equation}
  a_{\rm reh} = \left(\frac{g_{*S,0}}{g_{*S,\rm reh}}\right)^{1/3}
  \frac{T_0}{T_{\rm reh}}\,,
  \qquad
  a_i = a_{\rm reh}\,e^{-N}\,.
\end{equation}
Numerically, with $g_{*S,0} = 3.938$, $g_{*S,\rm reh} = 106.75$,
$T_0 = 2.7255\,\mathrm{K}$, $T_{\rm reh} = 10^{15}\,\mathrm{GeV}
= 1.160\times 10^{28}\,\mathrm{K}$, and $N = 60$:
\begin{align}
  \left(\frac{g_{*S,0}}{g_{*S,\rm reh}}\right)^{1/3}
  &= \left(\frac{3.938}{106.75}\right)^{1/3} = 0.3326\,, \\
  a_{\rm reh} &= 0.3326 \times \frac{2.7255}{1.160\times 10^{28}}
  = 7.813\times 10^{-29}\,, \\
  a_i &= 7.813\times 10^{-29}\times e^{-60}
  = 7.813\times 10^{-29}\times 8.757\times 10^{-27}
  = 6.842\times 10^{-55}\,.
\end{align}

\paragraph{Step 5: Comoving cutoff.}
\begin{equation}
  k_c = a_i\,k_{c,\,\rm phys}
  = a_i \times \frac{1}{l_P}\sqrt{\frac{\pi}{\ln D}}\,.
\end{equation}
To obtain $k_c$ in $\mathrm{Mpc}^{-1}$, we multiply by
$1\,\mathrm{Mpc} = 3.086\times 10^{22}\,\mathrm{m}$:
\begin{align}
  k_c &= 6.842\times 10^{-55} \times
  \frac{\sqrt{\pi/(2\times 10^{14})}}{1.616\times 10^{-35}\,\mathrm{m}}
  \times 3.086\times 10^{22}\,\mathrm{m}\,\mathrm{Mpc}^{-1} \nonumber\\
  &= 6.842\times 10^{-55} \times 2.455\times 10^{27}\,\mathrm{m}^{-1}
  \times 3.086\times 10^{22}\,\mathrm{m}\,\mathrm{Mpc}^{-1} \nonumber\\
  &\simeq 1.64\times 10^{-4}\,\mathrm{Mpc}^{-1}\,.
\end{align}
This confirms eq.~(12) in the main text.

\paragraph{Step 6: Tensor-to-scalar ratio.}
From $H_{\rm inf} = M_{\rm Pl}\sqrt{8\pi^2/\ln D}$ (where
$M_{\rm Pl} = m_{\rm Pl}/\sqrt{8\pi}$ is the reduced Planck mass):
\begin{equation}
  P_T = \frac{2}{\pi^2}\frac{H_{\rm inf}^2}{M_{\rm Pl}^2}
  = \frac{2}{\pi^2}\times\frac{8\pi^2}{\ln D}
  = \frac{16}{\ln D}\,,
\end{equation}
giving
\begin{equation}
  r = \frac{P_T}{A_s} = \frac{16}{A_s\,\ln D}\,.
\end{equation}
For $A_s = 2.1\times 10^{-9}$ and $\ln D = 2\times 10^{14}$:
$r = 16/(2.1\times 10^{-9}\times 2\times 10^{14}) = 3.81\times 10^{-5}$,
confirming eq.~(17).


% ============================================================
% APPENDIX B — MCMC Convergence Diagnostics
% ============================================================
\section{MCMC convergence diagnostics}
\label{app:convergence}

We present diagnostic tests for the MCMC analysis described in
section~\ref{sec:methodology}.

\paragraph{Chain statistics.}
The single Metropolis--Hastings chain comprises $24\,528$ samples (after 30\% burn-in removal) with a total weight of $82\,649$ (mean multiplicity~$3.4$), corresponding to an acceptance rate of $\sim 30$\% and an effective sample size $N_{\rm eff} \simeq 5\,400$ (computed as $N_{\rm eff} = (\sum w_i)^2 / \sum w_i^2$). The final Gelman--Rubin statistic, computed via the split-chain method, is $R - 1 < 0.005$, well below the conventional $R - 1 < 0.01$ target used in cosmological analyses.

\paragraph{Trace plots.}
Figure~\ref{fig:trace_plots} shows the trace plots and running weighted
means for the five key parameters: $\log_{10}(\ln D)$, $\alpha$, $H_0$,
$\tau$, and $n_s$. All running means stabilise within the first
$\sim 1\,000$ samples and show no systematic drift thereafter, indicating
adequate mixing. The CCR parameters $\log_{10}(\ln D)$ and $\alpha$
exhibit broader excursions than the $\Lambda$CDM parameters, reflecting
the flat posterior (the chain explores the full prior range uniformly).

\paragraph{Gelman--Rubin diagnostic.}
Figure~\ref{fig:gelman_rubin} shows the split-chain $R - 1$ statistic
as a function of sample number. All parameters fall well below the $R - 1 = 0.01$ target, with the final value of $R - 1 < 0.005$ confirming robust convergence.

\paragraph{$\chi^2$ trace.}
Figure~\ref{fig:chi2_trace} shows the $\chi^2$ trace, confirming that
the chain does not become trapped in local minima. The running mean stabilises at $\langle\chi^2\rangle = 1008.9$ with minimum $\chi^2_{\min} = 1002.6$.

\begin{figure}[p]
  \centering
  \includegraphics[width=\linewidth]{fig_trace_plots}
  \caption{Trace plots and running weighted means for the five key
    parameters. The CCR parameters explore the full prior range
    uniformly, while the $\Lambda$CDM parameters converge rapidly.}
  \label{fig:trace_plots}
\end{figure}

\begin{figure}[p]
  \centering
  \includegraphics[width=\linewidth]{fig_gelman_rubin}
  \caption{Split-chain Gelman--Rubin diagnostic $R - 1$ as a function
    of sample number. The dashed line marks the $R - 1 = 0.01$ target.
    All parameters fluctuate near or below $R - 1 = 0.01$ for the
    majority of the chain, with a final value of $R - 1 < 0.005$.}
  \label{fig:gelman_rubin}
\end{figure}

\begin{figure}[p]
  \centering
  \includegraphics[width=\linewidth]{fig_chi2_trace}
  \caption{$\chi^2$ trace plot. The running mean (black) stabilises at
    $\langle\chi^2\rangle = 1008.9$; the minimum is
    $\chi^2_{\min} = 1002.6$.}
  \label{fig:chi2_trace}
\end{figure}


% ============================================================
% APPENDIX C — Numerical Validation
% ============================================================
\section{Numerical validation}
\label{app:validation}

\paragraph{$\Lambda$CDM recovery.}
Figure~\ref{fig:lcdm_recovery} (upper panel) shows the fractional
difference $|\Delta C_\ell^{TT}/C_\ell^{TT}|$ between the
CCR-modified and standard $\Lambda$CDM spectra for increasing values
of $\ln D$. The results confirm the expected limiting behaviour:

\begin{center}
\begin{tabular}{ccc}
  \hline\hline
  $\ln D$ & $k_c$ [Mpc$^{-1}$] & $\max|\Delta C_\ell^{TT}/C_\ell^{TT}|$ \\
  \hline
  $10^{14}$ & $2.3\times 10^{-4}$ & $2.7\times 10^{-1}$ \\
  $10^{15}$ & $7.3\times 10^{-5}$ & $1.6\times 10^{-2}$ \\
  $10^{16}$ & $2.3\times 10^{-5}$ & $2.2\times 10^{-4}$ \\
  $10^{17}$ & $7.3\times 10^{-6}$ & $1.8\times 10^{-6}$ \\
  $10^{18}$ & $2.3\times 10^{-6}$ & $2.8\times 10^{-11}$ \\
  \hline\hline
\end{tabular}
\end{center}

\noindent
For $\ln D = 10^{16}$ (the upper edge of our prior range), the maximum
fractional deviation is $2.2\times 10^{-4}$, confirming that the
$\Lambda$CDM limit is recovered to high precision. For $\ln D \geq 10^{17}$,
the deviation drops below $10^{-5}$ across all multipoles.

\paragraph{$k$-grid convergence.}
Figure~\ref{fig:lcdm_recovery} (lower panel) compares the TT power
spectrum computed with $n_k = 500$, 1000, and 1500 $k$-grid points
against a reference computation with $n_k = 2000$, for
$\ln D = 2\times 10^{14}$ and $\alpha = 2$. The fractional differences
are below $10^{-5}$ for $n_k \geq 500$, confirming that our fiducial
choice of $n_k = 1500$ is more than adequate.

\begin{figure}[p]
  \centering
  \includegraphics[width=\linewidth]{fig_lcdm_recovery}
  \caption{\textbf{Upper:} $\Lambda$CDM recovery test. Fractional
    difference between CCR and $\Lambda$CDM $C_\ell^{TT}$ for
    increasing $\ln D$. The dashed line marks the
    $5\times 10^{-5}$ threshold.
    \textbf{Lower:} $k$-grid convergence test for
    $\ln D = 2\times 10^{14}$, comparing $n_k = 500$, 1000, 1500
    against a reference with $n_k = 2000$. All grids agree to better
    than $10^{-5}$.}
  \label{fig:lcdm_recovery}
\end{figure}